% This is text associated with the open textbook "Open Electromagnetics"
% See immediately below for author, date
% This work is licensed under the Creative Commons Attribution-ShareAlike 4.0 International License. (CC BY-SA 4.0) 
% To view a copy of the license, visit http://creativecommons.org/licenses/by-sa/4.0/

%ORB TITLE Comparison of Electrostatics and Magnetostatics 
%ORB AUTHOR 0        # S.W. Ellingson, Virginia Tech (ellingson@vt.edu)
%ORB DATE 20180615   # yyyymmdd
%ORB ABSTRACT Comparison of Electrostatics and Magnetostatics 

\label{m0115_Comparison_of_Electrostatics_and_Magnetostatics} % <-- Must match name of this file and module directory name.  Don't mess with this!

Students encountering magnetostatics 
for the first time have usually been exposed to electrostatics 
already.
Electrostatics and magnetostatics exhibit many similarities.
These are summarized in Table~\ref{m0115_tStaticsDuality}.
The elements of magnetostatics presented in this table are all formally introduced in other sections; the sole purpose of this table is to point out the similarities.  
%
\begin{table*}[t]
\begin{center}
\begin{tabular}{lll}
~ & {\bf electrostatics} & {\bf magnetostatics} \\
\hline\hline
Sources & static charge  & steady current, \\
~ &                      & magnetizable~material \\
\hline
Field intensity & ${\bf E}$ (V/m)  & ${\bf H}$ (A/m) \\
Flux density    & ${\bf D}$ (C/m$^2$) & ${\bf B}$ (Wb/m$^2$=T) \\
\hline
Material relations & ${\bf D}=\epsilon{\bf E}$ & ${\bf B}=\mu{\bf H}$ \\
~                  & ${\bf J}=\sigma{\bf E}$ &  \\
\hline
Force on charge $q$ & ${\bf F}=q{\bf E}$ & ${\bf F}=q{\bf v} \times {\bf B}$ \\
\hline
Maxwell's Eqs. & $\oint_{\mathcal S}{{\bf D}\cdot d{\bf s}} = Q_{encl}$ & $\oint_{\mathcal S}{{\bf B}\cdot d{\bf s}} = 0$ \\
~~(integral)  & $\oint_{\mathcal C}{{\bf E}\cdot d{\bf l}} = 0$ & $\oint_{\mathcal C}{{\bf H}\cdot d{\bf l}} = I_{encl}$ \\
\hline
Maxwell's Eqs. & $\nabla\cdot{\bf D}=\rho_v$ & $\nabla\cdot{\bf B}=0$ \\
~~(differential)   & $\nabla\times{\bf E}=0$ & $\nabla\times{\bf H}={\bf J}$ \\
\hline
Boundary Conditions & $\hat{\bf n} \times \left[{\bf E}_1-{\bf E}_2\right] = 0$     & $\hat{\bf n} \times \left[{\bf H}_1-{\bf H}_2\right] = {\bf J}_s$ \\
~~                  & $\hat{\bf n} \cdot \left[{\bf D}_1-{\bf D}_2\right] = \rho_s$ & $\hat{\bf n} \cdot \left[{\bf B}_1-{\bf B}_2\right] = 0$ \\
\hline
Energy storage & Capacitance & Inductance \\
Energy density & $w_e=\frac{1}{2}\epsilon E^2$ & $w_m=\frac{1}{2}\mu H^2$ \\
Energy dissipation & Resistance &  \\
\hline
\end{tabular}
\end{center}
\label{m0115_tStaticsDuality}
\caption{A summary of the duality between electrostatics and magnetostatics.}
\end{table*}

The technical term for these similarities is \emph{duality}.
\index{duality}
Duality also exists between voltage and current in electrical circuit theory.
For more about the concept of duality, see ``Additional Reading'' at the end of this section.

%%%%%%%%%%%%%%%%%%%%%%%%%%%%%%%%%%%%%%%%%

{\bf Additional Reading:}
\begin{itemize}
\item ``\href{https://en.wikipedia.org/wiki/Duality_(electricity_and_magnetism)}{Duality (electricity and magnetism)}'' on Wikipedia.
\item ``\href{https://en.wikipedia.org/wiki/Duality_(electrical_circuits)}{Duality (electrical circuits)}'' on Wikipedia.
\end{itemize}





